\documentclass[UTF8,a4paper,11pt]{ctexart}
\usepackage{geometry}
\usepackage{booktabs}
\usepackage{hyperref}
\usepackage{longtable}
\usepackage{xcolor}
\usepackage{tikz}
\usepackage{graphicx}
\geometry{left=2.6cm,right=2.6cm,top=2.6cm,bottom=2.6cm}
\hypersetup{colorlinks=true,linkcolor=black,citecolor=blue,urlcolor=blue}

\newcommand{\demofigure}[3]{
\begin{figure}[htbp]
\centering
\IfFileExists{figures/#1}{
\includegraphics[width=0.92\linewidth]{figures/#1}
}{
\fbox{\begin{minipage}[c][5.2cm][c]{0.88\linewidth}
\centering
\textbf{#2}\\[0.6em]
请将截图保存为 \texttt{report/figures/#1}
\end{minipage}}
}
\caption{#3}
\end{figure}
}

\newcommand{\IntentTrainSize}{59}
\newcommand{\IntentTestSize}{20}
\newcommand{\IntentAccuracy}{0.950}
\newcommand{\IntentMacroFone}{0.905}

\newcommand{\EvalQuestionCount}{50}
\newcommand{\EvalRouterAccuracy}{0.760}
\newcommand{\EvalSourceCoverage}{1.000}
\newcommand{\EvalAverageAnswerLength}{1184.3}


\title{基于证据约束检索增强生成的深圳北理莫斯科大学报考问答系统}
\author{NLP 大作业项目}
\date{2026年6月30日}

\begin{document}
\maketitle

\begin{abstract}
高校报考咨询是一类强事实、强时效且高度依赖个人条件的信息获取任务。通用生成模型虽然能够给出流畅文本，但在招生章程、专业证书、教学语言、招生人数和历年分数线等问题上容易产生无来源回答或过期记忆。本文构建了一个面向深圳北理莫斯科大学报考咨询的中文问答系统。系统以学校招生信息网、研究生招生信息网和公开招生资料为知识来源，将政策文本写入检索库，将分数线、专业层次、单双学籍、证书、教学语言、招生计划、学费和学制等信息写入结构化表，并通过轻量意图分类器决定检索、表格查询、联网搜索、澄清回复和日常对话的调用路径。系统在服务器端接入本地 Qwen2.5-0.5B-Instruct 模型，但将其限定为非结构化证据回答的语言组织模块；涉及分数、位次、证书和招生人数的事实问题由结构化查询和规则模板直接生成，以降低模型改写数字或虚构政策的风险。当前意图分类器使用 \IntentTrainSize 条训练样本和 \IntentTestSize 条测试样本，测试准确率为 \IntentAccuracy，宏平均 F1 为 \IntentMacroFone；系统评测集包含 \EvalQuestionCount 条问题，路由准确率为 \EvalRouterAccuracy，检索来源覆盖率为 \EvalSourceCoverage。实验和演示结果表明，证据约束的 RAG 工作流能够在课程项目规模内兼顾可运行性、可追溯性和事实边界。
\end{abstract}

\section{Introduction}

高考志愿填报是一类典型的高风险信息检索任务。考生和家长不仅需要了解学校和专业的基本情况，还需要将个人分数、位次、选科类别和专业偏好与历年录取数据联系起来。对于中外合作大学而言，培养模式、学籍归属、证书授予、教学语言和招生批次具有明显的学校特异性。若仅依赖普通搜索引擎，用户需要在多个网页之间自行整合信息；若直接依赖通用对话模型，回答又可能缺少可核验来源。

检索增强生成（retrieval-augmented generation, RAG）为这类问题提供了更合适的技术路线。RAG 的核心思想是在生成回答前先检索外部知识，并将回答约束在可追溯证据之内，从而降低语言模型仅依赖参数记忆时的事实错误风险\cite{lewis2020rag}。Qwen 系列模型具有中文、多语言、结构化输入输出和指令跟随能力；其中 Qwen2.5 提供从 0.5B 到 72B 的开源密集模型，并在预训练和后训练阶段扩展了数据规模和任务覆盖\cite{qwen2024technical,qwen25modelcard}。因此，本项目采用 RAG-first 设计：招生事实由数据层维护，Qwen 只承担受控生成，而不是作为事实数据库。

本文的目标不是构造泛化聊天机器人，而是实现一个面向深圳北理莫斯科大学报考咨询的可复现系统。系统重点处理学校信息、招生政策、专业介绍、历年分数线、研究生招生、校园生活和报考风险提示等问题。与将网页文本直接用于小模型微调不同，本文将事实知识保留在可更新、可引用的数据层中，并使用轻量训练模块完成问题类型识别和查询路由。这一设计使系统能够在信息更新后通过重新抓取和表格维护完成修订，而不需要重新训练生成模型参数。

\section{Method}

\subsection{任务定义}

给定用户中文问题 $q$ 和可选考生画像 $p$，其中 $p$ 包括省份、科类、分数、位次和意向专业，系统输出回答 $a$、问题类型 $t$、引用来源集合 $S$ 以及结构化记录 $R$。结构化记录覆盖普通高考分数线、本科专业单双学籍、证书来源、教学语言、招生人数、学费住宿费，以及硕士和博士专业的招生计划、学制和教学语言。系统的约束条件是：凡涉及分数、位次、证书、招生人数和政策条款的回答，必须能够追溯到结构化记录或来源片段；证据不足时返回澄清问题或缺失说明，而不是生成无来源结论。

\subsection{系统流程}

\begin{center}
\begin{tikzpicture}[node distance=1.4cm, every node/.style={draw, rounded corners, align=center, minimum width=3.2cm, minimum height=0.8cm}]
\node (src) {官方网页与表格};
\node (ingest) [below of=src] {抓取与清洗};
\node (store) [below of=ingest] {文档块与结构化招生表};
\node (route) [below of=store] {意图分类与路由};
\node (retrieve) [below left of=route, xshift=-1.7cm] {本地证据检索};
\node (score) [below right of=route, xshift=1.7cm] {分数线查询};
\node (web) [below of=retrieve, yshift=-0.4cm] {受控联网搜索};
\node (gen) [below of=route, yshift=-2.5cm] {证据约束回答生成};
\node (ui) [below of=gen] {Web Demo};
\draw[->] (src) -- (ingest);
\draw[->] (ingest) -- (store);
\draw[->] (store) -- (route);
\draw[->] (route) -- (retrieve);
\draw[->] (route) -- (score);
\draw[->] (route) -- (web);
\draw[->] (retrieve) -- (gen);
\draw[->] (web) -- (gen);
\draw[->] (score) -- (gen);
\draw[->] (gen) -- (ui);
\end{tikzpicture}
\end{center}

\subsection{数据层}

数据层分为文档证据和结构化事实两部分。文档证据来自本科招生信息网、招生章程、研究生招生简章、专业介绍、常见问题和学校公开页面。抓取脚本保存原始快照，将 HTML 正文清洗为纯文本，并切分为可检索片段；每个片段保留标题、链接、来源类型和可信度标记。结构化事实表用于保存需要精确查询的招生信息，字段包括年份、省份、科类、批次、录取类型、专业组、最低分、最低位次、专业层次、单双学籍、证书、教学语言、招生人数、学费、学制和来源链接。对于学校官网未以统一表格发布的历史投档数据，系统保留政府公开信息和第三方汇总来源标记，并在回答中提示最终以考试院和学校发布信息为准。

\subsection{检索、路由与回答生成}

系统采用先路由、后查询、再生成的流水线。首先，意图分类器和关键词规则将问题分为学校事实、招生政策、分数线查询、专业介绍、校园生活、奖助政策、就业深造、报考建议、专业结构化信息、研究生招生、日常对话、澄清问题和不支持问题。其次，路由结果决定查询路径：分数线和报考风险问题进入结构化分数线表；专业证书、教学语言和招生人数问题进入专业维度表；政策和学校事实问题进入文档检索；最新通知类问题在用户显式选择或问题触发时进入受控联网搜索。

当前检索模块实现了轻量词项匹配。系统将用户问题切分为中英文词项，对文档片段计算词项重叠分数，并返回标题、链接和摘要片段。联网搜索模块只面向深北莫招生网、研究生招生网和学校官网等官方域名；若联网失败，系统回退到本地资料库并在回答中写明边界。回答生成阶段将问题类型、结构化记录、检索片段和考生画像合并为中文回答。所有回答统一附带边界说明：系统只进行资料整理和风险提示，不能保证录取结果，不能替代学校或考试院的正式公告。

\subsection{轻量模型训练}

为体现可训练的 NLP 组件，本文训练了一个轻量意图分类器。分类器采用字符 n-gram TF--IDF 特征和逻辑回归模型，用于识别学校事实、招生政策、分数线查询、专业介绍、校园生活、奖助政策、就业深造、报考建议、日常对话和不支持问题等类别。该模块只参与查询路由，不承担招生事实记忆；当分类器不可用或置信度不足时，系统回退到关键词规则。对于“你是谁”“谢谢”“你会什么”“讲个笑话”等日常对话，系统不触发招生资料检索，也不交给生成模型自由发挥，而是使用受控模板以深北莫报考问答助手身份简短回复。

\subsection{本地 Qwen 生成模型与架构}

服务器端部署的生成模型为 Qwen2.5-0.5B-Instruct。该模型属于自回归因果语言模型，采用 decoder-only Transformer 架构\cite{vaswani2017attention,qwen2024technical,qwen25modelcard}。输入文本经分词器映射为 token 序列后进入词嵌入层，随后通过多层 Transformer decoder block 进行上下文建模。每个 decoder block 的核心包括带因果掩码的自注意力模块、前馈网络、残差连接和归一化层。因果掩码保证第 $i$ 个位置只能利用其左侧上下文，从而使模型能够按从左到右的方式生成回答。

在位置建模上，Qwen 使用旋转位置编码（rotary positional embedding, RoPE）将相对位置信息注入注意力计算\cite{su2024roformer}。在注意力结构上，本项目使用的 Qwen2.5-0.5B-Instruct 配置包含 24 层 decoder、14 个 query attention heads 和 2 个 key-value heads；这种 grouped-query attention（GQA）设计减少了推理时 key-value cache 的存储和读取开销，同时保留多查询头对上下文的表达能力\cite{ainslie2023gqa,qwen25modelcard}。在前馈网络中，Qwen 系列采用 SwiGLU 激活以增强非线性表示能力\cite{shazeer2020glu,qwen2024technical}；归一化采用 RMSNorm，以降低层归一化的计算复杂度并稳定深层 Transformer 训练\cite{zhang2019rmsnorm,qwen2024technical}。此外，Qwen2.5 模型使用 attention QKV bias、词嵌入共享和长上下文配置；0.5B Instruct 版本公开配置支持约 32k token 上下文窗口和约 8k token 生成长度\cite{qwen25modelcard}。

本项目没有对 Qwen 权重进行事实型微调。原因是招生章程、招生人数和分数线具有年度更新特征，将这些事实写入模型参数会增加过期记忆和幻觉风险。系统采用角色提示词调优和接口级约束：后端通过 OpenAI-compatible 接口调用本地 Qwen，并在 system prompt 中固定其身份为“深北莫报考问答助手”；对于非结构化证据类问题，Qwen 只根据检索片段组织语言；对于分数线、单双学籍、证书、教学语言和招生人数等结构化事实，后端跳过 Qwen 改写，直接返回表格查询结果。该设计把生成模型用于表达和归纳，把事实裁决保留在可审计的数据层中。

\section{Result}

\subsection{系统实现结果}

系统实现为一个基于 FastAPI 的网页演示系统。页面左侧记录省份、科类、分数、位次、意向专业和联网搜索开关，中间为对话区域，右侧展示本地检索来源、联网搜索结果、分数线表格、专业证书和研究生招生记录。该界面将回答文本和证据记录放在同一工作流中，便于检查系统是否真正进行了检索、联网搜索和结构化查询，而不是仅由生成模型直接输出回答。

\subsection{Demo 展示案例}

为展示系统的完整工作流，本文选取五类代表性问题作为结果截图。第一类是日常对话问题，用于验证系统能够识别“你是谁”“谢谢”“讲个笑话”等非招生检索问题，并以深北莫报考问答助手身份简短回复。第二类是广东物理类分数线查询，用于展示结构化最低分和最低位次记录。第三类是本科专业信息查询，用于展示单双学籍、证书、教学语言和招生人数等维度。第四类是研究生招生查询，用于展示硕士专业、教学语言、招生计划和证书信息。第五类是联网搜索场景，用于展示系统在最新通知类问题上优先返回官方域名来源。

\demofigure{demo-daily-chat.png}{日常对话与助手身份}{日常对话示例。用户输入“你是谁”或“讲个笑话”后，系统返回简短回答，不触发招生资料检索，也不输出无关来源。}

\demofigure{demo-score-query.png}{广东物理类分数线查询}{结构化分数线示例。用户填写广东、物理类并询问“广东物理类多少分能报”，系统返回历年最低分、最低位次和报考边界说明。}

\demofigure{demo-undergraduate-program.png}{本科专业证书与教学语言}{本科专业信息示例。用户询问“电子与计算机工程是单学籍还是双学籍”，系统返回专业、证书、教学语言和结构化来源。}

\demofigure{demo-graduate-admission.png}{研究生招生计划查询}{研究生招生示例。用户询问“纳米生物技术硕士是英语教学吗，招多少人”，系统返回教学语言、招生人数和莫斯科大学硕士学位证书信息。}

\demofigure{demo-web-search.png}{联网搜索官方通知}{联网搜索示例。用户勾选联网搜索并询问最新招生通知，系统优先展示深北莫官方域名结果，并在回答中标出联网搜索边界。}

\subsection{意图分类结果}

当前意图分类器的训练集规模为 \IntentTrainSize，测试集规模为 \IntentTestSize。测试准确率为 \IntentAccuracy，宏平均 F1 为 \IntentMacroFone。该结果说明系统具备可训练的查询路由模块；由于训练集规模有限，本文不将该结果解释为通用语义理解能力，而仅用于支持课程项目中的任务型意图识别。

\subsection{系统评测结果}

系统评测集覆盖学校事实、招生政策、分数线、本科专业证书、教学语言、研究生招生人数、校园生活、奖助政策、就业深造、报考建议、日常对话和拒答问题，共 \EvalQuestionCount 条。当前路由准确率为 \EvalRouterAccuracy，检索来源覆盖率为 \EvalSourceCoverage，平均回答长度为 \EvalAverageAnswerLength 字符。主要错误来源包括两类：一是单个问题同时包含多个意图，当前单标签分类器难以完整表达；二是官方资料未覆盖的主观体验问题只能返回资料边界说明，无法生成可验证结论。

\subsection{失败模式}

系统仍存在三方面局限。第一，当前分数线表仍需要持续补充官方历年录取数据；缺失字段不能由模型推断。第二，词项检索保证了实现简洁和可复现，但在同义表达、长问题和跨句检索上弱于向量检索与重排模型。第三，报考建议依赖考生位次、专业偏好、当年招生计划和志愿填报策略，系统只能提供资料整理和风险提示，不能给出确定录取结论。

\section{Conclusion}

本文实现了一个面向深圳北理莫斯科大学报考咨询的中文问答系统。系统将招生事实保留在可更新、可追溯的数据层中，通过检索和结构化查询约束回答，再用轻量训练模块完成意图路由，并在服务器端接入本地 Qwen2.5-0.5B-Instruct 模型用于非结构化证据回答的语言组织。该设计满足课程项目对数据抓取、模型组件、系统部署、网页演示和报告复现的要求，同时避免了小模型事实微调在招生场景下可能带来的过期记忆和不可解释回答。后续工作应继续扩展省份分数线覆盖，加入向量检索和重排模型，并扩大评测集以更稳定地衡量回答忠实性。

\begin{thebibliography}{99}
\bibitem{lewis2020rag}
Lewis, P. et al. Retrieval-Augmented Generation for Knowledge-Intensive NLP Tasks. \textit{Advances in Neural Information Processing Systems} (2020). \url{https://arxiv.org/abs/2005.11401}

\bibitem{qwen2024technical}
Qwen Team. Qwen2.5 Technical Report (2024). \url{https://arxiv.org/abs/2412.15115}

\bibitem{qwen25modelcard}
Qwen Team. Qwen2.5-0.5B-Instruct Model Card. Hugging Face (2024). \url{https://huggingface.co/Qwen/Qwen2.5-0.5B-Instruct}

\bibitem{qwen2025embedding}
Qwen Team. Qwen3 Embedding: Advancing Text Embedding and Reranking Through Foundation Models (2025). \url{https://arxiv.org/abs/2506.05176}

\bibitem{vaswani2017attention}
Vaswani, A. et al. Attention Is All You Need. \textit{Advances in Neural Information Processing Systems} (2017). \url{https://arxiv.org/abs/1706.03762}

\bibitem{su2024roformer}
Su, J. et al. RoFormer: Enhanced Transformer with Rotary Position Embedding. \textit{Neurocomputing} 568, 127063 (2024). \url{https://arxiv.org/abs/2104.09864}

\bibitem{ainslie2023gqa}
Ainslie, J. et al. GQA: Training Generalized Multi-Query Transformer Models from Multi-Head Checkpoints. \textit{Proceedings of EMNLP} (2023). \url{https://arxiv.org/abs/2305.13245}

\bibitem{shazeer2020glu}
Shazeer, N. GLU Variants Improve Transformer. arXiv (2020). \url{https://arxiv.org/abs/2002.05202}

\bibitem{zhang2019rmsnorm}
Zhang, B. and Sennrich, R. Root Mean Square Layer Normalization. \textit{Advances in Neural Information Processing Systems} (2019). \url{https://arxiv.org/abs/1910.07467}

\bibitem{smbu2025regulation}
深圳北理莫斯科大学招生信息网. 深圳北理莫斯科大学2025年夏季高考招生章程. \url{https://admission.smbu.edu.cn/info/1336/12222.htm}

\bibitem{smbuadmission}
深圳北理莫斯科大学招生信息网. \url{https://admission.smbu.edu.cn/}

\bibitem{szgov2025}
深圳市政府网站. 深圳多所高校公布2025年广东省本科批次普通类投档情况. \url{https://www.sz.gov.cn/cn/xxgk/zfxxgj/zwdt/content/post_12287337.html}

\bibitem{dxsbb2025}
大学生必备网. 2025深圳北理莫斯科大学录取分数线（含2023-2024历年）. \url{https://www.dxsbb.com/news/148133.html}
\end{thebibliography}

\end{document}
